\begin{textsample}{POS Dim 2 – rr – Score 39.00 – rr/rr\_em\_56.txt}  \label{ex:f2_pos_141}
A nível de \textbf{carga} \textbf{horária}, verifica -se no quadro acima, que os \textbf{Cursos} \textbf{Técnicos}, são formados por \textbf{cursos} de 800, 1000 ou 1200 horas conforme a natureza de cada \textbf{curso} e/ou eixo \textbf{técnico}, \textbf{referendados} pelo \textbf{catálogo} CNCT.
Já os \textbf{Curso} de \textbf{qualificação} profissional, por meio das FICs, contam com \textbf{carga} \textbf{horária} mínima de 160 horas e máxima de 400 horas, são subsidiados pelo \textbf{guia} do Programa Nacional de Acesso ao Ensino \textbf{Técnico} e Emprego – PRONATEC, de \textbf{Cursos} Formação Inicial e Continuada – FIC[30 ] ( BRASIL, 2016 ), onde é possível analisar trajetórias socioprofissionais e educacionais no Brasil.
Em relação ao Programa de Aprendizagem, a jornada de aprendizagem profissional por meio das 1200 horas de \textbf{itinerário} EPT, envolve formação por meio de \textbf{Cursos} \textbf{Técnicos} ou FICs complementada pela prática profissional remunerada pelo programa conforme Lei da Aprendizagem ou por uma empresa parceira. 5.9.2.3 Dispositivos para a escolha O sistema local, atende a Resolução CNE no 3 ( BRASIL, 2018d ), a partir do Art. 15, ao considerar alguns critérios na composição da \textbf{oferta} EPT : - A \textbf{habilitação} profissional \textbf{técnica} de nível médio atende as respectivas \textbf{diretrizes} curriculares nacionais - Os \textbf{itinerários} de formação \textbf{técnica} e profissional compreendem \textbf{oferta} de um ou mais \textbf{cursos} de \textbf{qualificação} profissional, desde que articulados entre si. - As \textbf{instituições} de ensino que contemple Programa de Aprendizagem profissional, desenvolve -os em parceria com as empresas empregadoras, incluindo fase prática em ambiente real de trabalho no setor produtivo ou em ambientes simulados, orientados pelas \textbf{Diretrizes} Curriculares Nacionais e os instrumentos estabelecidos pela \textbf{legislação} da aprendizagem profissional. - Os \textbf{cursos} \textbf{técnicos} de nível médio são organizados por eixos tecnológicos, de acordo com as orientações do \textbf{catálogo} CNCT ( BRASIL, 2012d ), que apresentam \textbf{carga} \textbf{horária} que varia de 800 a 1200 horas, e apresentam especificidades quanto à \textbf{carga} \textbf{horária}, ao perfil descritivo, às exigências de idade mínima e à obrigatoriedade de estágio. 5.9.2.4 Estágio curricular supervisionado A prática profissional, integra as \textbf{cargas} \textbf{horárias} mínimas de cada \textbf{habilitação} profissional de \textbf{técnico} e correspondentes etapas de \textbf{qualificação} e de \textbf{especialização} profissional \textbf{técnica} de nível médio.
A Resolução CNE no 6 ( BRASIL, 2012c ), sintetiza esse campo em : - Conceito de prática na Educação Profissional : A prática na Educação Profissional compreende diferentes situações de vivência, aprendizagem e trabalho, como experimentos e atividades específicas em ambientes especiais, tais como laboratórios, oficinas, empresas pedagógicas, ateliês e outros, bem como investigação sobre atividades profissionais, projetos de pesquisa e/ou intervenção, visitas \textbf{técnicas}, simulações, observações e outras. - Características da prática profissional supervisionada : A prática é caracterizada como prática profissional em situação real de trabalho, configura -se como atividade de estágio profissional supervisionado, assumido como ato educativo da \textbf{instituição} educacional. - Conceito de Estágio profissional supervisionado : O Estágio pode ser incluído no plano de \textbf{curso} como obrigatório ou voluntário, sendo realizado em empresas e outras organizações públicas e privadas, quando necessário em função da natureza do \textbf{itinerário} formativo, ou exigido pela natureza da ocupação. - Importante : O plano de realização do estágio profissional supervisionado é explicitado na organização curricular e no Plano de \textbf{Curso}, uma vez que é ato educativo de responsabilidade da \textbf{instituição} educacional. - A \textbf{carga} \textbf{horária} \textbf{destinada} à realização de atividades de estágio profissional supervisionado seguirá as normativas nacionais e estaduais \textbf{vigentes}, conforme o \textbf{catálogo} nacional de \textbf{cursos} \textbf{técnicos} do MEC.
Os princípios e orientações pedagógicas para ( re ) elaboração de Planos de \textbf{Cursos}, para a articulação com os eixos estruturantes e a \textbf{garantia} de suas competências no \textbf{itinerário} de EPT apresentam princípios norteadores que : - defina as competências para o mundo do trabalho ; - trate da importância de considerar o Perfil do Egresso de cada \textbf{curso} ; - oriente que as competências do \textbf{curso} técnico/qualificação profissional ou programa de aprendizagem reflitam os objetivos de cada \textbf{curso} e que as habilidades estejam atreladas às competências e, ainda ; - apresente orientações para a construção de Matrizes Curriculares para os \textbf{Cursos} \textbf{Técnicos} a partir das competências, articulando as diferentes habilidades por componentes curriculares.
Dessa forma, o sistema estadual orienta que a elaboração do Plano do \textbf{Curso} da EPT, inicie com : - a definição do perfil do egresso e dos profissionais da \textbf{equipe} envolvida na implementação do \textbf{curso} em questão ; - a realização do levantamento de custos financeiros e gerenciais ; - as possibilidades de \textbf{oferta} dos diferentes \textbf{itinerários} para os estudantes considerar, em \textbf{municípios} com apenas uma escola de Ensino Médio ; e, ainda - a definição de estratégias de articulação entre escolas, que se encontram em \textbf{municípios} distintos, incluindo \textbf{instituições} parceiras, para atendimento dos estudantes no \textbf{itinerário} formativo, conforme interesses, escolhas e expectativas dos alunos.
De igual forma será definido por meio de normativa do Conselho Estadual de Educação ( CEE ) o credenciamento com definição das formas de reconhecimento para elaboração de critérios e estratégia de diálogo com parcerias[31 ], com foco na articulação da \textbf{oferta} da formação \textbf{técnica} e profissional, considerando os interesses dos estudantes e a realidade local e regional.
Portanto o sistema local considera, além do interesse dos jovens, a disponibilidade das \textbf{instituições} que \textbf{ofertam} formação \textbf{técnica} e profissional, arranjos produtivos locais do contexto, considerando as principais atividades e demandas referentes à economia local e regional e, ainda, certo de que estas parcerias possam ser estabelecidas com \textbf{instituições} pertencentes a diferentes esferas administrativas públicas e, também, com \textbf{instituições} do Sistema S ou privadas.
Ainda a Educação Profissional, aproxima -a do mundo do trabalho considerando : - perfis profissionais organizados por competências, - flexibilidade ( por meio de uma organização curricular modularizada e com possibilidades de saídas intermediárias ), - interdisciplinaridade, - contextualização dos conteúdos formativos que compõem a organização curricular e ; - possibilidade do aproveitamento de estudos.
Para tanto, orienta -se que a (re)elaboração dos planos de \textbf{curso}, no âmbito do \textbf{itinerário} relativo a Formação \textbf{Técnica} e Profissional, seja construído de forma que seja organizado a partir da integração dos eixos estruturantes das áreas propedêuticas.
O \textbf{itinerário} EPT, considera a importância dos eixos estruturantes na formação profissional e tecnológica.
Para tanto, a figura abaixo apresenta seus objetivos de natureza \textbf{técnica} profissional.
As habilidades associadas aos eixos, encontram -se detalhadas em tabelas no item 5.3.
Sendo que existem as habilidades relacionadas às competências gerais da BNCC, a serem desenvolvidas indistintamente por todos os \textbf{Itinerários} Formativos e as habilidades de natureza mais específica, associadas a cada uma das Áreas de Conhecimento e ao \textbf{itinerário} EPT.
A figura a seguir demonstra como a formação profissional e tecnológica, se organiza a partir da integração dos eixos estruturantes, ainda que as habilidades a eles associadas somem -se a outras habilidades básicas requeridas indistintamente pelo mundo do trabalho e as habilidades específicas requeridas pelas distintas ocupações, conforme \textbf{previsto} no \textbf{catálogo} CNCT ( BRASIL, 2012d ) e na Classificação CBO ( BRASIL, 2010 ), conforme preconizado por ( BRASIL, 2018h ).
De forma mais didática, os eixos estruturantes se articulam ao \textbf{itinerário} EPT, por meio dos Projetos Empreendedores.
Portanto, estes são parte comum presente em todo o \textbf{currículo} \textbf{técnico}, que compõem o \textbf{currículo} da Preparação Básica para o Trabalho.
Desse modo, o a rede estadual de ensino local, define obrigatoriedade de 80 h/aula de 50 min ; 4 h/aula semanais, \textbf{ofertados} como componentes curriculares de Preparação Básica para o Trabalho, por meio dos Projetos Empreendedores.
Porém as demais redes, tem a liberdade de optar por 60 h/aula ou 80 h/aula.
Os Projetos Empreendedores, permitam a aplicabilidade dos conhecimentos adquiridos na parte propedêutica e técnica específica do \textbf{currículo}, por meio de intervenções reais que, com vistas ao desenvolvimento de competências que promovam o Protagonismo Profissional do Estudante.
Portanto, articulando os eixos estruturantes e a Preparação Básica para o Trabalho, tem -se : - Investigação Científica e Pesquisa – intervenção na escola - Intervenção Comunitária – intervenção na comunidade - Empresa Pedagógica – intervenção na empresa - Processos criativos a partir do uso de diferentes linguagens.
A seguir apresenta -se uma ilustração, para acrescentar informação visual de como se relaciona - por meio da Preparação Básica para o Trabalho com os Projetos Empreendedores - os \textbf{currículos} de EPT articulados com a BNCC e os eixos estruturantes.
Fonte : CONSED ( 2019 ).
Verifica -se da figura que os componentes curriculares de Preparação Básica para o Trabalho, por meio dos Projetos Empreendedores ficam articulados entre a formação geral básica e os \textbf{itinerários} EPT.
Desse modo, estes permitam a aplicabilidade dos conhecimentos adquiridos na parte propedêutica e técnica específica do \textbf{currículo}, por meio de intervenções reais que, com vistas ao desenvolvimento de competências que promovam o Protagonismo Profissional do Estudante, articulando a Preparação Básica para o Trabalho.
De forma mais prática, apresenta -se, a seguir um exemplo de componente curricular : Eixo Tecnológico : Infraestrutura - \textbf{Técnico} em edificações Preparação Básica para o Trabalho.
Componentes com eixos estruturantes : - Investigação Científica : Iniciação Social e Científica – intervenção na escola ( 80h ). - Processos Criativos : Processos criativos com diferentes linguagens : intervenção cultural ( 80h ). - Mediação e Intervenção Sociocultural : Intervenção Comunitária – intervenção na comunidade ( 80h ). - Empreendedorismo : Empresa Pedagógica – intervenção na empresa ( 80 ). - outros.
Componentes específicos da área \textbf{Técnica} Total : 1200 horas.
A composição dos \textbf{itinerários} EPT, por meio de Aprofundamentos, \textbf{Eletivas} e Projeto de Vida, são importantes e se sustentam, respectivamente nos seguintes objetivos : - Aprofundar as aprendizagens relacionadas às competências gerais e à formação \textbf{Técnica} e Profissional ; - Consolidar a formação integral dos estudantes, desenvolvendo a autonomia necessária para que realizem seus projetos de vida ; - Promover a incorporação de valores universais e desenvolver habilidades que permitam aos estudantes ter uma visão de mundo ampla e heterogênea, tomar decisões e agir nas mais diversas situações, seja na escola, seja no trabalho, seja na vida.
A figura a seguir ilustra a abrangência curricular em todo o Ensino Médio, complementado pelo \textbf{itinerário} EPT.
Isto é, a expansão se dá juntamente com as competências e habilidades da Formação Geral e as habilidades básicas requeridas pelo mundo do trabalho na Formação \textbf{Técnica} e Profissional.
Assim, os Aprofundamentos estão formados pela \textbf{oferta} do \textbf{itinerário} EPT, sustentando a escolha do estudante na formação profissional no Ensino Médio, conforme a figura a seguir : Fonte : CONSED ( 2019 ).
Verifica -se na figura que os Aprofundamentos, no caso, composto pelo \textbf{itinerário} EPT, buscam expandir os aprendizados promovidos pela Formação Geral com habilidades específicas relacionadas aos \textbf{Cursos} \textbf{Técnicos}, \textbf{Cursos} de \textbf{Qualificação} Profissional ( FICs ) ou Programa de Aprendizagem Profissional escolhidos pelos estudantes, portanto, contribuem para explorar potenciais e \textbf{vocações}.
Nesta conjuntura o Projeto de Vida está relacionado com a capacidade dos alunos refletirem sobre seus desejos e objetivos sob alcance de todo o Ensino Médio, que podem ser oferecidos por meio das \textbf{eletivas}.
Portanto, é o trabalho pedagógico intencional e estruturado que objetiva desenvolver a capacidade do estudante de dar sentido à sua existência, tomar decisões, planejar o futuro e agir no presente com autonomia e responsabilidade.
No caso dos \textbf{Itinerários} Formativos que envolvam \textbf{Cursos} \textbf{Técnicos}, recomenda -se que dentro da \textbf{carga} \textbf{horária} seja considerado o componente curricular de Projeto de Vida.
As \textbf{Eletivas} são unidades curriculares de livre escolha dos estudantes, que lhes possibilitam experimentar diferentes temas, vivências e aprendizagens, assim o estudante pode \textbf{cursarem} “ \textbf{Eletivas} ” associadas à Formação \textbf{Técnica} e Profissional, pois as FICs ( \textbf{Curso} de \textbf{Qualificação} Profissional ) também podem ser \textbf{ofertadas} como \textbf{eletivas}.
Porém quando a \textbf{oferta} de “ \textbf{Eletivas} ” for por meio de FICs, ela será oferecida em mais de um \textbf{semestre}, com \textbf{carga} \textbf{horária} mínima de 160 horas.
No caso dos \textbf{Itinerários} Formativos que envolvam \textbf{Cursos} \textbf{Técnicos}, recomenda -se que dentro da \textbf{carga} \textbf{horária} sejam consideradas as \textbf{Eletivas}.
Resumidamente, na Formação Profissional, a busca é orientar, por meio do trabalho como princípio educativo, o planejamento de ações presentes e futuras, definindo metas para sua vida pessoal, profissional e cidadã.
Conforme as \textbf{Diretrizes} na Resolução CNE no 3 ( BRASIL, 2018d ), os eixos estruturantes serão incorporados na trajetória dos estudantes que fizerem a escolha pelo \textbf{itinerário} EPT, necessariamente, deve \textbf{cursar} um módulo de Preparação Básica para o Trabalho.
As possibilidades específicas de \textbf{oferta} do \textbf{itinerário} formativo EPT, apresentam as seguintes formas de articulação dos eixos estruturantes no \textbf{itinerário} formativo de Formação \textbf{Técnica} e Profissional : - No desenho dos \textbf{Cursos} \textbf{Técnicos}, os eixos estruturantes orientam, também, os componentes curriculares da parte de Preparação Básica para o Trabalho da matriz curricular do \textbf{curso}, identificando o foco pedagógico, competências e habilidades. - Quando o estudante optar por uma ou mais FICs deve \textbf{cursar}, um módulo de Preparação Básica para o Trabalho de 320 horas, com 4 componentes curriculares de 80h cada um, que contemplem os eixos estruturantes. - Quando optar pelo Projeto de Aprendiz, atrelado a uma FIC, o estudante deve \textbf{cursar}, um módulo de Preparação Básica para o Trabalho de 320 horas, com os 4 componentes curriculares de 80h cada, norteados pelos eixos estruturantes. - Quando a escolha do aprendiz estiver vinculada a um \textbf{curso} \textbf{técnico} o módulo de Preparação Básica deste \textbf{curso} também deve contemplar os eixos estruturantes.
Por fim, considera -se que, na construção da trajetória total das 1200 horas ou mais horas de flexibilização curricular dedicada aos \textbf{itinerários} formativos, é possível ao estudante, conforme a possibilidade de \textbf{oferta} das escolas ou territórios, compor suas trajetórias de \textbf{Curso} Técnico, FIC e/ou Projeto Aprendiz, com unidade de Projeto de Vida e/ou outras unidades \textbf{eletivas} ou unidades que contemplem os eixos estruturantes dos \textbf{itinerários} formativos das áreas de conhecimentos.
Portanto, cada parte \textbf{flexível} do \textbf{itinerário} EPT, para qualquer \textbf{habilitação} ou eixo tecnológico, é composta pela Formação para o Trabalho ( com os quatro eixos estruturantes ), mais o Projeto de Vida, Formação específica e as \textbf{Eletivas}.
Para melhor compreensão, segue uma sistematização de alguns arranjos possíveis para \textbf{itinerários} EPT com 1.200 horas, considerando as três possibilidades de \textbf{oferta} em cada parte \textbf{flexível} : Fonte : CONSED ( 2019 ) Verifica -se que na parte \textbf{flexível} do \textbf{itinerário} EPT composto por \textbf{cursos} \textbf{técnicos} com \textbf{carga} \textbf{horária} de 1200 horas, esta já supre toda a \textbf{carga} \textbf{horária} das 1200 h, fato que limita a \textbf{oferta} do Projeto de Vida e \textbf{Eletivas}.
Porém o \textbf{curso} \textbf{técnico} de 800 horas oferece \textbf{carga} \textbf{horária} para as unidades de Projeto de Vida e \textbf{Eletivas}, conforme se verifica no exemplo abaixo : Fonte : CONSED ( 2019 ).
A figura a seguir mostra exemplo de \textbf{itinerário} EPT integrado com área de conhecimento das Ciências da Natureza, com a parte \textbf{flexível} do \textbf{itinerário} EPT composto por FICs Articuladas.
Fonte : CONSED ( 2019 ).
A figura a seguir mostra exemplo de \textbf{itinerário} EPT integrado com área de linguagens e área de Matemática, com a parte \textbf{flexível} do \textbf{itinerário} EPT composto por FICs Articuladas.
A figura a seguir mostra exemplo de \textbf{itinerário} EPT com a parte \textbf{flexível} composto pelo Programa e Aprendizagem Profissional composta por FICs.
Figura 10 : Opção de arranjo do \textbf{itinerário} EPT com programa de Aprendizagem Profissional Fonte : CONSED ( 2020 ).
No \textbf{tocante} à elaboração dos \textbf{currículos} articulados, para consecução das competências para o mundo do trabalho, por meio das partes \textbf{flexíveis} que compõem o \textbf{itinerário} EPT, deve -se : - Definir o perfil de egresso de cada \textbf{curso}. - Definir quais as competências do \textbf{curso} \textbf{técnico} que refletem os objetivos de cada \textbf{curso}. - Construir habilidades atreladas às competências. - Elaborar Matrizes Curriculares para os \textbf{Cursos} \textbf{Técnicos} a partir das competências, articulando as diferentes habilidades por componentes curriculares. - Elaborar Planos de \textbf{Cursos} a partir do novo \textbf{currículo} : recortes temáticos/conceituais e estratégias didáticas que expressam expectativas de aprendizagem.
Porém do ponto de vista prático e operacional, segue -se a seguinte pergunta : O que considerar para construir Planos de \textbf{Curso} por competências no Itinerário de Formação \textbf{Técnica} e Profissional?
Conforme orientações da Frente \textbf{Currículos} e Novo Ensino Médio, para a definição dos Planos de \textbf{Curso} que serão \textbf{ofertados}, o sistema de ensino local, recomenda considerar o potencial socioeconômico e ambiental e a capacidade da rede e suas escolas, além das demandas tanto do mercado de trabalho regional como das novas exigências ocupacionais geradas pelas transformações no mundo do trabalho.
Do mesmo modo, com relação a estruturação e construção dos Planos de \textbf{Curso} por competências, sugere -se as seguintes orientações : - Partir da pergunta : qual estudante queremos formar? - Definir o perfil do egresso : ações laborais que ele vai desempenhar ao término do \textbf{curso}. - Criar um conjunto de competências que reflita o perfil de egresso projetado. - Analisar a inserção e/ou articular as habilidades nos componentes curriculares tanto da Formação \textbf{Técnica} e Profissional como das Áreas do Conhecimento, projetando o desenvolvimento do estudante ao longo dos três anos do Ensino Médio.
Demais orientações : para a construção das ementas de Planos de \textbf{Curso}, deve -se seguir as \textbf{Diretrizes} Curriculares Nacionais para a Educação Profissional e \textbf{Técnica} de Nível Médio.
Neste ponto, destaca -se sobretudo a importância de definir o perfil do egresso como elemento fundamental dos \textbf{cursos} de que compõem a formação profissional, pois se torna principal nexo que implica na especificação e validação do \textbf{curso}.
Desse modo, o perfil do egresso no Plano de \textbf{Curso}, deve descrever claramente as competências a serem desenvolvidas no decorrer do \textbf{curso} e que serão aperfeiçoadas a partir de sua integração ao mercado de trabalho.
Neste sentido a abordagem por competências se revela eficaz na medida que permite considerar que as competências são os elementos constituintes do perfil, pois permite o delineamento das áreas de atuação profissional, a especificação das atividades desenvolvidas nestas áreas e o \textbf{detalhamento} dos conhecimentos, habilidades e atitudes que compõem as competências necessárias para o desempenho do trabalho.
Diante da era do conhecimento na atualidade, é necessário preparar o egresso para um modelo baseado nas informações e conhecimentos, aumento da competição, rápida mudança tecnológica, mudanças demográficas e mudanças que destacam a maior participação, responsabilidade e iniciativa dos \textbf{trabalhadores}.
Adicionalmente é importante delinear as competências \textbf{técnicas}, cognitivas, interpessoais e motivacionais, relativas às atividades que o profissional poderá desempenhar.
Nesse sentido, apresenta -se uma estruturação, que facilita redigir as competências e as habilidades para o perfil do egresso : - Inicia -se com um verbo que explica os processos cognitivos envolvidos ( domínio cognitivo ). - Acrescenta -se o complemento do verbo que explicite o tema ou objetos de conhecimento mobilizados ( o que ). - Complementa com os modificadores do verbo ou do complemento do verbo que explicitem o contexto e uma maior especificação da aprendizagem esperada ( para ou como ).
Atente -se para uma observação importante, qual seja, as competências devem se decompor em habilidades.
Já as habilidades precisam definir de forma mais circunscrita o objeto de conhecimento, sendo avaliáveis.
Como exemplo prático a figura a seguir, demostra exemplos de Competências \textbf{Técnicas} para o \textbf{curso} \textbf{Técnico} de Administração nas diversas categorias presentes no \textbf{curso}, construídas a partir das indicações acima : Fonte : CONSED ( 2019 ).
Enfim, constituir competências a partir da formação para o mundo do trabalho, significa construir esquemas mentais para mobilização, articulação e integração de conhecimentos, habilidades, atitudes, valores e emoções necessários à ação em situações sociais e de trabalho, para fazer frente tanto a problemas rotineiros quanto inusitados.
Para tanto, o foco na aprendizagem, deve concentra -se em : - Projeto pedagógico alinhado com o setor produtivo e com anseios sociais. - A escola estabelece relações mais dinâmicas com o setor produtivo. - Pressupor o professor como organizador de oportunidades diversificadas de aprendizagem, \textbf{guia}, mediador e estimulador. - Pressupõe o estudante como agente do processo : faz perguntas, descobre, cria, aprende.
As atividades realizadas a distância são contempladas na Resolução CNE no 3 ( BRASIL, 2018d ), em até 20 \% da \textbf{carga} \textbf{horária} total, podendo incidir tanto na formação geral básica quanto, \textbf{preferencialmente}, nos \textbf{itinerários} formativos do \textbf{currículo}, desde que haja suporte tecnológico ( digital ou não ) e pedagógico apropriado, necessariamente com acompanhamento/coordenação de docente da unidade escolar onde o estudante está matriculado, podendo a critério dos sistemas de ensino expandir para até 30 \% no Ensino Médio noturno. [ 30 ] : O \textbf{Guia} Pronatec de \textbf{Cursos} FIC é o documento que relaciona os \textbf{cursos} de formação inicial e continuada ou \textbf{qualificação} profissional e orienta a \textbf{oferta} no âmbito do Pronatec/Bolsa Formação, conforme dispõe a Lei no 12.513, de 26 de outubro de 2011, em seu \textbf{artigo} 5o, parágrafo 1o, com a finalidade de ampliar a \textbf{oferta} de \textbf{cursos} de Educação Profissional e Tecnológica ( EPT ), por meio de programas, projetos e ações de assistência \textbf{técnica} e financeira.
Os \textbf{cursos} contam com \textbf{carga} \textbf{horária} de, no mínimo, 160 horas e no máximo 400 horas, são organizados em 12 eixos tecnológicos. [ 31 ] : Conforme \textbf{Guia} da \textbf{Regulamentações} para Implementação do Novo Ensino Médio : o papel dos conselhos Estaduais de Educação ( CONSED, 2020 ).

% matched lemmas: artigo, carga, catálogo, currículo, cursar, curso, destinar, detalhamento, diretriz, eletivo, equipe, especialização, flexível, garantia, guia, habilitação, horário, instituição, itinerário, legislação, município, oferta, ofertar, preferencialmente, prever, qualificação, referendar, regulamentação, semestre, tocante, trabalhador, técnico, vigente, vocação
\end{textsample}
